\chapter{总结}

\section{工作总结}

本作品围绕“如何让 Agent 的真实行动处于可控状态”这一问题，设计并实现了 Aegis Runtime。系统不把安全理解为模型在生成文本时的一次自我判断，而是把风险判断、策略与权限、人工审批、任务图调度、Effect 生命周期、提交与回滚、Audit 记录放入连续的运行时过程。用户提交的任务因此不再只是一个黑盒执行过程：系统能够说明某个节点为何等待、某项资源是否已经提交、取消后哪些影响被恢复、哪些独立成果被保留。

作品以 Transactional Controlled Execution、Approval-aware TaskGraph Scheduling 和 Dependency-aware Selective Rollback 为三项核心创新。第一项机制使可恢复副作用具有 Pending、校验、提交或回滚的明确边界，避免将候选结果过早视为可信成果；第二项机制让审批状态参与任务图调度，使等待审批的节点及其受影响范围得到约束，同时允许安全独立节点继续推进；第三项机制依据任务、请求、Effect、Checkpoint 和资源目标之间的关联进行局部恢复，避免因为一个问题误删全部独立成果。

三组正式实验分别检验安全处置、图调度和状态恢复，浏览器场景进一步呈现 Grant、Deny、Cancel、危险操作阻断和 Audit 记录等用户可见闭环。实验和场景共同说明，作品的价值不只在于“检测到风险”，而在于风险判断能够改变真实执行路径；不只在于“支持审批”，而在于审批决定能够影响图的推进；不只在于“可以回滚”，而在于回滚能够按影响范围处理并保护独立成果。系统最终验收通过。

\section{作品意义}

对使用者而言，Aegis Runtime 将复杂的 Agent 执行过程转化为可以理解的状态：等待审批不再是模糊的停顿，提交不再只是没有依据的成功提示，取消也不再意味着不确定的后台行为。用户能够在任务、审批、Effect 和 Audit 视图之间追踪同一事项，并据此作出授权、拒绝或取消决定。脱敏展示使这种可见性不以暴露密钥、原始不可信内容或隐藏推理为代价。

对开发者而言，统一的关联关系为系统测试和问题排查提供了基础。风险判断、审批事件、图节点、Effect 和资源处理不再各自留下难以对应的日志，而是围绕 task、request 和 Checkpoint 等事实建立联系。这样，当出现异常时，可以检查问题发生在风险判定、审批处理、调度推进还是资源恢复阶段；当需要增加新的可恢复工具时，也可以沿用既有的受控执行与 Effect 记录方式。

从更广泛的 Agent 安全实践看，本作品说明运行时治理需要同时考虑预防、约束和恢复。只做输入过滤无法覆盖工具调用之后的状态变化，只做人工审批无法保证审批决定与任务调度一致，只做异常捕获也无法决定哪些成果应被保留。将三类问题放在同一条可追溯链路中，能够使安全机制从单点规则进一步发展为可验证的系统行为。

\section{局限性}

本系统是单进程、受控环境中的 Runtime 原型。当前原型未覆盖分布式工作流编排、服务重启恢复、完整 ACID 语义、通用 MCP/Skill 协议和同一节点 speculative execution。选择性回滚仅适用于已经纳入 Effect、Checkpoint 和关联关系管理的可恢复状态；对于真实支付、邮件发送或其他不可逆外部动作，系统不能宣称具有通用的事后撤销能力，只能在动作释放前通过审批、策略和提交边界降低风险。

实验采用固定 fixtures、确定性规则和受控时延，重点验证机制行为、状态一致性和结果可复核性。Safety 正式结论不依赖 LLM 裁决，Graph 实验也不用于说明统计显著性或真实大模型在开放环境中的性能。开放网络中的输入变化、第三方服务差异、长时间运行与多进程协调都会带来额外挑战，因此本作品的结论应理解为对当前工程范围内机制的验证，而不是对所有 Agent 应用场景的普适保证。

\section{未来工作}

后续工作可以在不改变安全边界的前提下，逐步扩展状态协调能力。例如，可研究 TaskGraph、审批与 Effect 状态在受控持久化条件下的恢复方式，并明确恢复过程中应当重新核验的授权和资源条件；对于数据库写入、网络请求和第三方 API，可进一步区分可逆、可补偿和不可逆操作，并为不同类别设计更明确的提交条件与用户确认方式。

在交互层面，未来可以继续优化风险理由、审批摘要和 Audit 查询，使不同技术背景的用户更容易理解系统为何提出确认、为何阻断某个动作以及为何保留某项独立成果。在实验层面，可在保持可复核和不夸大结论的前提下，增加更多类型的受控业务流程，并考察不同资源冲突与长期任务场景下的状态一致性。无论功能如何扩展，核心原则仍应保持：任何新增执行能力都需要先说明其授权边界、可观察状态和异常后的处理方式。
